\part{Das Urteil}

\chapter{Der zweite Entwurf im Großen}\label{kap:zweiter}

\begin{merksatz}
\boxtitel{Was hier zusätzlich zählt}
\begin{center}
{\small
\renewcommand{\arraystretch}{1.5}
\begin{tabular}{@{}P{3.2cm}P{2.2cm}P{2.6cm}@{}}
\toprule
\textbf{Kriterium} & \textbf{Band 20} & \textbf{Hier}\\
\midrule
Verstehenszeit & \textcolor{gruen}{ja} & \textcolor{gruen}{ja}\\
Änderbarkeit & \textcolor{gruen}{ja} & \textcolor{gruen}{ja}\\
Prüfbarkeit & \textcolor{gruen}{ja} & \textcolor{gruen}{ja}\\
Größenordnung & \textcolor{gruen}{ja} & \textcolor{gruen}{ja}\\
\midrule
\textbf{Übersteht Teilausfall} & --- & \textcolor{gruen}{\textbf{ja}}\\
\textbf{Im Betrieb befragbar} & --- & \textcolor{gruen}{\textbf{ja}}\\
\textbf{Einzeln erneuerbar} & --- & \textcolor{gruen}{\textbf{ja}}\\
\bottomrule
\end{tabular}}
\end{center}

\begin{center}
\begin{minipage}{0.88\textwidth}
\kursiv{Sieben Bedingungen, alle prüfbar --- und die drei unteren lassen sich nicht nachträglich herstellen. Sie folgen aus der Bauform, und wer sie beim ersten Entwurf übergeht, baut beim zweiten das Werk um, nicht den Code.}
\end{minipage}
\end{center}
\end{merksatz}

\begin{merksatz}
\boxtitel{Der Fall: die Verlaufsansicht durch das Werk}
\textbf{Der erste Entwurf.}

\begin{center}
{\small
\renewcommand{\arraystretch}{1.45}
\begin{tabular}{@{}P{3.2cm}P{2.4cm}P{2.4cm}@{}}
\toprule
\textbf{Schwäche} & \textbf{Kapitel} & \textbf{Wirkung}\\
\midrule
Nur ein Zeitpunkt & \ref{kap:zeit} & \textcolor{akzent}{\textbf{Verlauf falsch}}\\
Einheit implizit & \ref{kap:zahlen} & \textcolor{akzent}{\textbf{Faktor $18$}}\\
$n{+}1$ über die Naht & \ref{kap:zeit_hin} & \textcolor{akzent}{$5{,}1$ s}\\
Keine Zeitschranken & \ref{kap:fehler} & \textcolor{akzent}{\textbf{Kaskade}}\\
Fehler verschwiegen & \ref{kap:fehler} & \textcolor{akzent}{\textbf{Lücke unsichtbar}}\\
Keine Verfolgung & \ref{kap:beobachtbarkeit} & \textcolor{akzent}{nicht aufklärbar}\\
Wiederholung ohne Kennung & \ref{kap:naht} & \textcolor{akzent}{Werte doppelt}\\
\bottomrule
\end{tabular}}
\end{center}

\begin{center}
\begin{minipage}{0.88\textwidth}
\kursiv{Sieben Schwächen, und jede stammt aus einem anderen Kapitel. Keine davon ist ein Codefehler --- jede ist eine Entscheidung, die an einer Naht nicht getroffen wurde. Der Code selbst kann in jedem Bestandteil vorbildlich sein.}
\end{minipage}
\end{center}

\textbf{Der zweite Entwurf.}

\begin{center}
{\small
\renewcommand{\arraystretch}{1.45}
\begin{tabular}{@{}P{2.8cm}P{2.0cm}P{2.0cm}P{2.0cm}@{}}
\toprule
& \textbf{Erster} & \textbf{Zweiter} & \\
\midrule
Zeilen gesamt & $4200$ & $\mathbf{5100}$ & $+21\,\%$\\
Felder je Befund & $6$ & $\mathbf{11}$ & \\
$95.$ Perzentil & $410$ ms & $\mathbf{95}$ ms & $4{,}3\times$\\
Bei Langzeitpatient & $\mathbf{5{,}1}$ s & $\mathbf{298}$ ms & $17\times$\\
Teilausfall verkraftet & \textcolor{akzent}{nein} & \textcolor{gruen}{\textbf{ja}} & \\
Aufklärbar & \textcolor{akzent}{nein} & \textcolor{gruen}{ja} & \\
Schwächen & $\mathbf{7}$ & $\mathbf{0}$ & \\
\bottomrule
\end{tabular}}
\end{center}

\begin{center}
\kursiv{Einundzwanzig Prozent mehr Zeilen, fünf zusätzliche Felder je Befund --- und sieben Fehlerklassen fort. Das Muster aus Band 20 wiederholt sich: Der zweite Entwurf ist länger.}
\end{center}

\begin{technik}
\boxtitel{Was die fünf Felder kosten}
\begin{center}
\renewcommand{\arraystretch}{1.45}
\begin{tabular}{@{}P{2.8cm}P{2.2cm}P{2.6cm}@{}}
\toprule
\textbf{Größe} & \textbf{Zuwachs} & \\
\midrule
Byte je Befund & $16 \to \mathbf{34}$ & $+112\,\%$\\
Bei $10^9$ Befunden & \textcolor{akzent}{$+\mathbf{18}$ \textbf{GB}} & \textcolor{mittel}{tragbar}\\
Übertragungs\-dauer & $+3\,\%$ & \textcolor{gruen}{unerheblich}\\
\textbf{Beantwortbare Fragen} & \textcolor{gruen}{\textbf{$+5$}} & \textcolor{gruen}{\textbf{der Gewinn}}\\
\bottomrule
\end{tabular}
\end{center}

\begin{center}
\begin{minipage}{0.86\textwidth}
\kursiv{Achtzehn Gigabyte für eine Milliarde Befunde --- bei heutigen Preisen ein zweistelliger Eurobetrag im Jahr. Dagegen stehen fünf Fragen, die vorher unbeantwortbar waren: Welcher Zeitpunkt? Welche Einheit? Wer hat wann was gesehen? Wurde korrigiert? Ist der Wert vollständig?}
\end{minipage}
\end{center}
\end{technik}
\end{merksatz}

\begin{aufgabe}
\begin{enumerate}[label=(\alph*)]
\item Prüfen Sie die sieben Bedingungen.
\item Bestimmen Sie die Reihenfolge des Umbaus.
\item Untersuchen Sie, was nicht mehr zu ändern war.
\end{enumerate}

\loesung
\textbf{(a) Die sieben Bedingungen.}

\begin{center}
{\small
\renewcommand{\arraystretch}{1.45}
\begin{tabular}{@{}P{3.0cm}P{2.4cm}P{2.6cm}@{}}
\toprule
\textbf{Bedingung} & \textbf{Prüfung} & \textbf{Ergebnis}\\
\midrule
Verstehenszeit & $5$ Fragen & \textcolor{mittel}{$88$ s}\\
Änderbarkeit & neues Feld & \textcolor{gruen}{ohne Bruch}\\
Prüfbarkeit & je Bestandteil & \textcolor{gruen}{\textbf{einzeln}}\\
Größenordnung & $1000$ Geräte & \textcolor{gruen}{$495$ k/s}\\
Teilausfall & \textbf{fünf Versuche} & \textcolor{gruen}{\textbf{bestanden}}\\
Befragbar & Verfolgung & \textcolor{gruen}{ja}\\
Einzeln erneuerbar & \textbf{Rollentausch} & \textcolor{gruen}{ja}\\
\bottomrule
\end{tabular}}
\end{center}

\begin{center}
\begin{minipage}{0.88\textwidth}
\kursiv{Die erste Zeile ist wieder der Preis: achtundachtzig Sekunden gegen einundsechzig beim ersten Entwurf. Ein Werk mit Zeitschranken, Wiederholungen und Verfolgung ist schwerer zu lesen als eines ohne --- so wie in Band 21 und 22.}
\end{minipage}
\end{center}

\begin{center}
\kursiv{Die Antwort ist dieselbe wie dort und keine Rücknahme: Zerlegen und einen Kommentar an die Naht, der das Zusammenspiel erklärt. Danach liegt die Verstehenszeit bei zweiundsechzig Sekunden.}
\end{center}

\medskip
\textbf{(b) Die Reihenfolge des Umbaus.}

\begin{center}
{\small
\renewcommand{\arraystretch}{1.45}
\begin{tabular}{@{}P{1.0cm}P{3.2cm}P{2.0cm}P{2.2cm}@{}}
\toprule
& \textbf{Schritt} & \textbf{Aufwand} & \textbf{Wirkung}\\
\midrule
$1$ & \textbf{Verfolgung einbauen} & $4$ PT & \textcolor{gruen}{\textbf{sehen, was ist}}\\
$2$ & Zeitschranken & $2$ PT & \textcolor{gruen}{keine Kaskaden}\\
$3$ & \textbf{Vertrag erweitern} & $6$ PT & \textcolor{gruen}{\textbf{Zeit, Einheit}}\\
$4$ & Beide Seiten umstellen & \textcolor{akzent}{$14$ PT} & \textbf{Kap.~\ref{kap:vertraege}}\\
$5$ & Aufrufe bündeln & $2$ PT & \textcolor{gruen}{$17\times$}\\
$6$ & Fehler sichtbar machen & $3$ PT & \\
\bottomrule
\end{tabular}}
\end{center}

\begin{center}
\begin{minipage}{0.88\textwidth}
\kursiv{Schritt eins steht vorn, obwohl er nichts verbessert: Ohne Verfolgung weiß man nicht, ob die späteren Schritte wirken. Das ist dieselbe Regel wie in Band 21 und 22 --- messen vor ändern.}
\end{minipage}
\end{center}

\begin{center}
\kursiv{Schritt vier ist der teuerste und der langwierigste: fünf Schritte über Wochen, weil beide Seiten nicht gleichzeitig umgestellt werden können. Das ist der Preis der Naht, den Kapitel~\ref{kap:vertraege} beziffert hat.}
\end{center}

\medskip
\textbf{(c) Was nicht mehr zu ändern war.}

\begin{center}
{\small
\renewcommand{\arraystretch}{1.5}
\begin{tabular}{@{}P{3.0cm}P{2.2cm}P{2.6cm}@{}}
\toprule
\textbf{Gegenstand} & \textbf{Änderbar?} & \\
\midrule
Code eines Bestandteils & \textcolor{gruen}{\textbf{leicht}} & \\
Naht zwischen zweien & \textcolor{mittel}{mühsam} & \textbf{Wochen}\\
\textbf{Format nach außen} & \textcolor{akzent}{\textbf{kaum}} & \textcolor{akzent}{\textbf{zwölf Fremdsysteme}}\\
Bedeutung eines Feldes & \textcolor{akzent}{nie} & \textcolor{akzent}{\textbf{Altbestand}}\\
Kennungs\-schema & \textcolor{akzent}{nie} & \\
\bottomrule
\end{tabular}}
\end{center}

\begin{center}
\begin{minipage}{0.88\textwidth}
\kursiv{Die vierte Zeile ist die härteste Grenze des Umbaus: Der gespeicherte Zeitpunkt war über Jahre der Speicherzeitpunkt. Man kann das Feld künftig richtig füllen --- die Millionen alten Werte bleiben, wie sie sind, und der Verlauf bleibt für sie falsch.}
\end{minipage}
\end{center}

\begin{center}
\renewcommand{\arraystretch}{1.45}
\begin{tabular}{@{}P{3.0cm}P{2.2cm}P{2.8cm}@{}}
\toprule
\textbf{Umgang mit Altbestand} & \textbf{Tauglich?} & \\
\midrule
Nachträglich ableiten & \textcolor{akzent}{\textbf{nein}} & \textcolor{akzent}{Information fehlt}\\
Schätzen & \textcolor{akzent}{\textbf{nie}} & \textcolor{akzent}{\textbf{erfundene Daten}}\\
\textbf{Als unbekannt kennzeichnen} & \textcolor{gruen}{\textbf{ja}} & \textcolor{gruen}{\textbf{ehrlich}}\\
Stichtag einführen & \textcolor{gruen}{ja} & \textbf{ab da verlässlich}\\
\bottomrule
\end{tabular}
\end{center}

\begin{center}
\kursiv{Die zweite Zeile ist die Versuchung, der man widerstehen muss: Aus dem Speicherzeitpunkt eine Abnahmezeit zu schätzen erzeugt Daten, die richtig aussehen und es nicht sind. Die dritte Zeile ist unbequem und die einzige zulässige.}
\end{center}

\begin{warnung}
\boxtitelr{Die Naht überlebt alles}
\begin{center}
\renewcommand{\arraystretch}{1.45}
\begin{tabular}{@{}P{2.8cm}P{2.2cm}P{2.6cm}@{}}
\toprule
\textbf{Gegenstand} & \textbf{Lebensdauer} & \\
\midrule
Eine Zeile Code & $\approx2$ Jahre & \\
Ein Bestandteil & $\approx7$ Jahre & \\
Eine Sprache im Werk & $\approx12$ Jahre & \\
\textbf{Ein Datenformat} & \textcolor{akzent}{$\approx\mathbf{20}$ \textbf{Jahre}} & \\
\textbf{Ein Kennungs\-schema} & \textcolor{akzent}{\textbf{unbegrenzt}} & \textcolor{akzent}{\textbf{Patientendaten}}\\
\bottomrule
\end{tabular}
\end{center}

\begin{center}
\begin{minipage}{0.88\textwidth}
\kursiv{Die letzte Zeile ist die längste Verpflichtung, die man in einem Klinikvorhaben eingeht: Eine Patientenkennung muss die Aufbewahrungsfristen überdauern --- dreißig Jahre und mehr. Sie wird jeden Bestandteil, jede Sprache und jeden Beteiligten überleben.}
\end{minipage}
\end{center}

\begin{center}
\kursiv{Deshalb rechtfertigt eine Naht einen Aufwand, der für Code unangemessen wäre. Zwei Wochen für ein Format, das zwanzig Jahre gilt, sind gut angelegt --- und zwei Tage sind es nicht.}
\end{center}
\end{warnung}
\end{aufgabe}
\vspace{-4pt}
\noindent Fowler\zit{8} liefert die Schritte des Um\-baus im Großen; Ford und Richards\zit{9} ordnen die Ent\-schei\-dun\-gen ein; Newman\zit{17} behandelt die schritt\-wei\-se Ver\-än\-de\-rung ge\-wach\-se\-ner Werke.
